\section{Questão 5}

\subsection{Enunciado}

Considere o seguinte conjunto de dados:
\begin{verbatim}
data = {
    'Idade': [28, 35, 45, 22, 31],
    'Produtos_Comprados': [5, 9, 7, 3, 6]
}
\end{verbatim}
Aplique uma técnica de transformação linear para criar novas features que correlacionem a Idade com Produtos\_Comprados.

\subsection{Solução}

A aula lista a criação de features de interação entre as transformações lineares. A operação é o produto de duas colunas, no mesmo formato de \texttt{valor\_total = preço \(\times\) quantidade}, e serve para correlacionar dois atributos numa coluna só. A feature nova é \texttt{Idade\_x\_Produtos} = Idade \(\times\) Produtos comprados. A correlação de Pearson entre as colunas originais resume a associação linear que esse produto carrega.

\subsection{Resultado}

A correlação de Pearson entre idade e produtos comprados é 0,729: na amostra, idade maior acompanha mais compras. A Tabela~\ref{tab:q5-interacao} traz o produto em cada linha. O arquivo \path{scripts/results/q5_interacao.csv} guarda a tabela.

\begin{table}[H]
    \centering
    \caption{Feature de interação entre idade e produtos comprados.}
    \label{tab:q5-interacao}
    \footnotesize
    \setlength{\tabcolsep}{8pt}
    \begin{tabular}{@{}rrr@{}}
        \toprule
        Idade & Produtos comprados & Idade \(\times\) produtos \\
        \midrule
        28 & 5 & 140 \\
        35 & 9 & 315 \\
        45 & 7 & 315 \\
        22 & 3 &  66 \\
        31 & 6 & 186 \\
        \bottomrule
    \end{tabular}
\end{table}
